\documentclass[12pt]{article}
\usepackage{graphicx}
\usepackage{amsmath,amssymb}
\usepackage{hyperref} %
\hypersetup{colorlinks=true, urlcolor=blue, linkcolor=blue, citecolor=red} %

\usepackage{lipsum} % for dummy text only
\graphicspath{{Figures/}}

\title{Testing Physics}
\author{FirstName LastName\\
Lab Partner(s): Name(s)\\
Physics 206 S20}
\date{\today}
%\keywords{keyword1, keyword2, keyword3, etc.}

\begin{document}
\maketitle

\begin{abstract}
% Abstract text here. 
\lipsum[1] % \lipsum creates filler text with similar number of frequency of letters as you would see in typical English writing. Delete and replace with your own content.
\end{abstract}

\section{Introduction}
% Text for Introduction. Include theory if needed. Referencing the lab write-up is a good way to keep the report concise.

\lipsum[2]
% Here is how you enter an equation
\begin{equation} \label{eq:RestMass}
E = mc^2
\end{equation}
\lipsum[3]

\section{Methods}
\subsection{Apparatus}
% Describe your apparatus. 
\lipsum[4]

% This is how you include a figure. Note the directory "figs/"
\begin{figure}  
    \centering
    \includegraphics[width=\linewidth]{figs/GSFC_20171208_Archive_e001988}
    \caption{Fermi bubble.\cite{fermibubble}}
    \label{fig:my_label}
\end{figure}

\subsection{Procedure}
% Procedure text here
\lipsum[5]

\section{Results}
% Results and figures below. 
\lipsum[6]

\section{Discussion}
\subsection{Analysis}
\lipsum[7]  
Something to be cited.\cite{goeppert_mayer}
\lipsum[2]

\subsection{Validation}
\lipsum[9]

\section{Conclusion}
\lipsum[10]

\section*{Acknowledgements}
I thank the A for working with me as a lab partner. I thank B for letting me use their equipment and providing access to various research tools, and I thank C and D for discussion and comments on this manuscript.

\begin{thebibliography}{}
% see Overleaf help for more info: https://www.overleaf.com/learn/latex/Bibliography_management_with_bibtex

\bibitem{fermibubble}
\href{images-assets.nasa.gov/image/GSFC_20171208_Archive_e001988/GSFC_20171208_Archive_e001988~orig.jpg}{"NASA Image ID: GSFC\_20171208\_Archive\_e001988"}

\bibitem{goeppert_mayer} 
Maria Goeppert Mayer. 
\textit{Nuclear Configurations in the Spin-Orbit Coupling Model. I. Empirical Evidence}.
Phys. Rev., 78:16, 1950.

\end{thebibliography}

\end{document}
